% Generated draft: requires mathematical and source review.
\documentclass[11pt]{article}
\usepackage[margin=25mm]{geometry}
\usepackage{fontspec}
\setmainfont{Latin Modern Roman}
\usepackage{amsmath,amssymb,mathtools,unicode-math}
\setmathfont{Latin Modern Math}
\usepackage{enumitem,needspace,xurl,hyperref}
\setlength{\parindent}{0pt}
\setlength{\parskip}{5pt}
\setlength{\emergencystretch}{3em}
\newcommand{\N}{\mathbb{N}}
\newcommand{\Z}{\mathbb{Z}}
\newcommand{\Q}{\mathbb{Q}}
\newcommand{\R}{\mathbb{R}}
\newcommand{\C}{\mathbb{C}}
\newcommand{\F}{\mathbb{F}}
\newcommand{\A}{\mathbb{A}}
\newcommand{\PP}{\mathbb{P}}
\newcommand{\E}{\mathbb{E}}
\newcommand{\eps}{\varepsilon}
\newcommand{\dd}{\,\mathrm{d}}
\begin{document}
\section*{UnsolvedMath ID: 30006241}
\textbf{Source problem code:} catalogue entry
\subsection{Definitions and mathematical statement}
Under the usual finite permutation-group interpretation, let \(\Omega\) be a finite nonempty set and let \(G\leq\operatorname{Sym}(\Omega)\). The action is transitive if, for every \(x,y\in\Omega\), some \(g\in G\) satisfies \(g(x)=y\). A finite group \(G\) is almost simple if there is a finite nonabelian simple group \(S\) such that
\[
S\leq G\leq\operatorname{Aut}(S),
\]
where \(S\) is identified with its group of inner automorphisms.

The action is binary if, for every integer \(n\geq2\) and every two tuples \((x_1,\ldots,x_n),(y_1,\ldots,y_n)\in\Omega^n\), the implication
\[
\left[\forall\,1\leq i<j\leq n\ \exists\,g_{ij}\in G:\quad g_{ij}(x_i)=y_i\ \text{and}\ g_{ij}(x_j)=y_j\right]
\Longrightarrow
\left[\exists\,g\in G\ \forall\,1\leq i\leq n:\quad g(x_i)=y_i\right]
\]
holds. Repeated tuple entries are allowed. A transitive action is primitive if its only \(G\)-invariant partitions are the partition into singletons and the partition with one part.

\textbf{Classification problem.} Classify all finite transitive binary permutation groups \(G\leq\operatorname{Sym}(\Omega)\) for which \(G\) is almost simple, without assuming primitivity. For this draft, classification is taken up to permutation isomorphism: two actions \(G\curvearrowright\Omega\) and \(H\curvearrowright\Lambda\) are identified when there are a group isomorphism \(\alpha:G\to H\) and a bijection \(f:\Omega\to\Lambda\) such that
\[
f(g(x))=\alpha(g)(f(x))\qquad(g\in G,\ x\in\Omega).
\]

\textbf{Unresolved scope qualifications.} The finite setting is suggested by the report title and context, but the extracted question does not explicitly state finiteness. Writing \(G\leq\operatorname{Sym}(\Omega)\) imposes faithfulness; the source does not explicitly say whether actions with nontrivial kernel are also intended. The equivalence relation for classification is likewise unstated. The formulation above therefore records a conventional interpretation requiring confirmation against the report, rather than resolving these omissions silently.

The supplied background describes this as an extension of an existing classification of primitive binary actions and mentions partial classifications for alternating and symmetric groups. The original passage refers to an unstated theorem formerly conjectured by Cherlin; its precise statement is absent from the supplied record and is not reconstructed here.
\subsection{Short English statement}
Classify the transitive binary actions of almost simple groups, allowing imprimitive actions. Binary means that whenever every corresponding pair of entries of two tuples can be matched by a group element, one group element matches the entire tuples. This draft uses finite faithful actions up to permutation isomorphism; these scope conventions require source confirmation. The supplied open status and classification background are dated curation, not independently verified conclusions.
\subsection{Sources}
\begin{itemize}
\item Local registry: \texttt{data/unsolvedmath/problems.json}, permanent UnsolvedMath ID \(30006241\), source code \texttt{catalogue entry}; title: \emph{Transitive Binary Actions of Almost Simple Groups}. The record supplies proposed year 2025 but identifies no individual proposer.
\item Source report: \emph{Finite Groups, Fusion Systems and Applications} (2025), Oberwolfach Reports, pp.\ 689--752, DOI \texttt{10.4171/owr/2025/16}. Supplied report link, also listed under the literature-source label ``current classification status'': \url{https://doi.org/10.4171/owr/2025/16}.
\item Supplied literature source: ``Gill--Guillot, alternating binary actions'': \url{https://arxiv.org/abs/2303.06003}.
\item The original passage attributes an unspecified theorem, formerly a conjecture of Cherlin, to work by various authors including Cherlin, Dalla Volta, Gill, Hunt, Liebeck, Spiga and Wiscons, with report reference labels [2, 3, 5, 6]. The theorem statement and the corresponding bibliography entries are not supplied.
\item Literature assessment checked at \texttt{2026-08-22}: the supplied curation labels the problem open, reports a classification of primitive binary actions and explicit partial classifications in the alternating and symmetric cases, and describes the general transitive almost simple case as a substantial open task. This is supplied dated curation derived from the Oberwolfach record, not an independently verified literature assessment; its status and cited claims require independent verification.
\item The registry labels the statement \texttt{corrected\_verified}, explaining that only formula delimiters, notation, grammar and typography were repaired while preserving mathematical content and scope. This is a supplied editorial qualification, not mathematical verification of this draft.
\end{itemize}
\end{document}
